\section{Problem statement}
As part of the Growing Roots research project\footnote{Growing Roots. From \url{https://growingroots.nl}, retrieved at 16-11-2024}, a device called the ``virtual window" is developed by the company TexTown Games\footnote{Design2Connect - TexTown Games. From \url{https://textowngames.nl/portfolio/design2connect/}, retrieved at 10-1-2025}. This device tracks the position of a person in front of a monitor and projects a landscape on the screen based on where the person is standing. Currently, only handcrafted landscapes are used for testing the system. For further research related to emotional impressions of landscapes, more environments are needed that differ in visually distinguishable parameters, such as the density of trees or the color and types of vegetation used. Handcrafting landscapes takes a significant amount of time, so the ability to create more terrains automatically is required. The landscapes generated will be landscapes of nature, with some man-made structures such as roads or benches. 

The generated terrain will only be used for the virtual window, which means that the camera will not move around the landscape. This limits the size and amount of terrain that needs to be generated. Therefore, each feature's order and properties will be generated with this rule in consideration.

Research in procedural terrain generation has often focused on generating large-scale terrain, but the virtual window only allows a smaller view on the terrain. This results in the generated terrain having a larger focus on smaller details such as vegetation and visible parts of the terrain.

As further research will focus on the effect of different terrain features, the researcher needs control over the terrain they want to create. This control should allow the user to select broader features of the terrain, such as the kind of ecology the terrain represents, and finer features, such as the density of vegetation or ruggedness of the ground.

As the system will be used by people who are not necessarily designers or programmers, the generation of the terrain will be automated with little user input. Because of this, we must be able to verify whether the generated terrain adheres to required constraints. These features can be broad for the entire terrain, such as that some desired biomes are included, or specific to the point of view, such that a large part of the generated is visible from the viewpoint.